%% O Brasil como Phi-engine Condicional: Estado Logístico e Coerência Hegemônica
%% Formatado para: RBPI (Revista Brasileira de Política Internacional)
%% Ou: Contexto Internacional (PUC-Rio)
\documentclass[12pt]{article}
\usepackage[utf8]{inputenc}
\usepackage[T1]{fontenc}
\usepackage[brazil]{babel}
\usepackage{amsmath,amssymb,amsthm,mathtools}
\usepackage{geometry}
\geometry{a4paper, margin=2.5cm}
\usepackage{setspace}\doublespacing
\usepackage{booktabs, array, longtable}
\usepackage{graphicx}
\usepackage[hidelinks]{hyperref}
\usepackage{natbib}
\bibliographystyle{plainnat}
\usepackage{lineno}
\linenumbers

\theoremstyle{definition}
\newtheorem{definition}{Defini\c{c}\~{a}o}
\newtheorem{proposition}{Proposi\c{c}\~{a}o}
\theoremstyle{remark}
\newtheorem{remark}{Observa\c{c}\~{a}o}

\providecommand{\tightlist}{\setlength{\itemsep}{0pt}\setlength{\parskip}{0pt}}

\hypersetup{
  pdftitle={O Brasil como Phi-engine Condicional},
  pdfauthor={Luis Fernando Marcelino},
}

\begin{document}

\begin{center}
{\LARGE\bfseries O Brasil como \(\Phi\)-\textit{engine} condicional:\\
Estado Log\'{i}stico, coer\^{e}ncia hegem\^{o}nica\\
e a transi\c{c}\~{a}o de fase 2028--2031}

\vspace{1.2em}
{\large Luis Fernando Marcelino}\\[0.3em]
{\normalsize Pesquisador Independente, Mil\~{a}o, It\'{a}lia}\\[0.3em]
{\small\texttt{lfrmrc@gmail.com}}

\vspace{0.5em}
{\small \textit{Agradecimentos:} V-Dem Institute, Banco Mundial.}
\end{center}

\vspace{1.5em}

% ─── RESUMO ──────────────────────────────────────────────────────────────────
\begin{abstract}
Este artigo aplica o modelo DEQ\textsuperscript{2} (\textit{Discursivo-Egem\^{o}nico Qu\^{a}ntico} ao quadrado) ao Brasil do per\'{i}odo 2024--2026, argumentando que o pa\'{i}s constitui o \'{u}nico ator do sistema global contempor\^{a}neo com a estrutura formal de um produtor de coer\^{e}ncia exportável (\(\Phi\)-\textit{engine}): n\~{a}o por dota\c{c}\~{a}o material superior, mas pelo assento operativo do Estado Log\'{i}stico (Cervo, 2008) associado \`{a} dupla presid\^{e}ncia BRICS--COP30 de 2025, \`{a} capacidade de media\c{c}\~{a}o estrutural documentada, e ao U-turn democr\'{a}tico verificado pelo V-Dem. O argumento formal \'{e} derivado do funcional de estabilidade sistêmica \(T_s^{\text{sys}}\), que exibe acoplamento assim\'{e}trico entre antifragilidade (\(A\), ampliada por \(\Phi\)) e dissipa\c{c}\~{a}o involutiva (\(\Omega\), dividida por \(\Phi\)). Demonstra-se que o Brasil preenche as cinco condi\c{c}\~{o}es definitórias de \(\Phi\)-\textit{engine} no sentido DEQ\textsuperscript{2}, mas que a fun\c{c}\~{a}o \'{e} \textbf{condicional} ao resultado das elei\c{c}\~{o}es presidenciais de 4 de outubro de 2026 --- o teste de falsificabilidade mais estrito do modelo no horizonte 2028--2031. S\~{a}o apresentadas quatro predi\c{c}\~{o}es datadas e falsific\'{a}veis pr\'{e}-registradas publicamente.

\medskip
\noindent\textbf{Palavras-chave:} Brasil; DEQ\textsuperscript{2}; Estado Log\'{i}stico; Kuramoto; transi\c{c}\~{a}o hegem\^{o}nica; BRICS; COP30; eleições 2026

\medskip
\noindent\textbf{Recebido:} [data de submiss\~{a}o]. \textbf{Aceito:} [data de aceita\c{c}\~{a}o].
\end{abstract}

\clearpage

% ─── 1. INTRODUÇÃO ───────────────────────────────────────────────────────────
\section{Introdu\c{c}\~{a}o: a lacuna entre BRICS e bifurca\c{c}\~{a}o}

O Brasil de 2025--2026 ocupa uma posi\c{c}\~{a}o estruturalmente an\^{o}mala no sistema internacional. Presidiu simultaneamente o BRICS (17\textdegree{} summit, Rio de Janeiro, julho de 2025) e a COP30 (Bel\'{e}m, novembro de 2025), produzindo em ambos documentos operativos n\~{a}o triviais. Manteve rela\c{c}\~{o}es diplom\'{a}ticas ativas com EUA, China, R\'{u}ssia, Ir\~{a} e Israel simultaneamente, sem ader\^{e}ncia a nenhum campo ideol\'{o}gico-hegemônico exclusivo. Registrou no V-Dem um \textit{U-turn} democr\'{a}tico no bienio 2022--2025. A Suprema Corte Federal condenou um ex-presidente por tentativa de golpe de Estado --- a primeira vez na hist\'{o}ria brasileira.

E, apesar disso, a aprova\c{c}\~{a}o do governo Lula em abril de 2026 est\'{a} no n\'{i}vel hist\'{o}rico m\'{i}nimo, com o candidato da oposi\c{c}\~{a}o radical estatisticamente empatado num segundo turno hipot\'{e}tico.

Esta combina\c{c}\~{a}o --- for\c{c}a estrutural externa e fragilidade pol\'{i}tica interna --- n\~{a}o \'{e} anedótica. \'{E} a assinatura emp\'{i}rica de um sistema em bifurca\c{c}\~{a}o expl\'{i}cita: duas trajet\'{o}rias poss\'{i}veis, uma data de revers\~{a}o conhecida, um mecanismo formal mensur\'{a}vel.

Este artigo situa o Brasil neste quadro atrav\'{e}s do modelo DEQ\textsuperscript{2} \citep{marcelino2026c}, derivado da Topografia discursiva \citep{marcelino2026a} e da Geometria da Hegemonia \citep{marcelino2026b}. O argumento central \'{e}: o Brasil \'{e} um \(\Phi\)-\textit{engine} condicional --- produtor de coer\^{e}ncia hegem\^{o}nica exportável se e somente se a elei\c{c}\~{a}o de outubro de 2026 mantenha a trajet\'{o}ria do Estado Log\'{i}stico operativa.

\subsection{Estrutura do artigo}

A Se\c{c}\~{a}o~2 apresenta o aparato DEQ\textsuperscript{2} m\'{i}nimo necess\'{a}rio. A Se\c{c}\~{a}o~3 define formalmente o \(\Phi\)-\textit{engine} e demonstra que o Brasil preenche as cinco condi\c{c}\~{o}es. A Se\c{c}\~{a}o~4 analisa o Estado Log\'{i}stico (Cervo) como realiza\c{c}\~{a}o pol\'{i}tica dessa estrutura. A Se\c{c}\~{a}o~5 examina as tr\^{e}s evidências de 2025: BRICS, COP30, condena Bolsonaro. A Se\c{c}\~{a}o~6 analisa as tr\^{e}s fragilidades reais. A Se\c{c}\~{a}o~7 apresenta os tr\^{e}s cen\'{a}rios eleitorais e suas consequ\^{e}ncias para \(T_s^{\text{sys}}\). A Se\c{c}\~{a}o~8 apresenta quatro predi\c{c}\~{o}es datadas pr\'{e}-registradas. A Se\c{c}\~{a}o~9 conclui.

% ─── 2. APARATO FORMAL ───────────────────────────────────────────────────────
\section{O aparato DEQ\textsuperscript{2}: elementos essenciais}

\subsection{O parâmetro de ordem de Kuramoto}

O modelo de Kuramoto \citep{kuramoto1984} descreve uma popula\c{c}\~{a}o de \(N\) osciladores com fases \(\theta_j(t)\) e frequ\^{e}ncias naturais \(\omega_j\) acoplados por
\[
\frac{d\theta_j}{dt} = \omega_j + \frac{K}{N}\sum_{k=1}^{N}\sin(\theta_k - \theta_j).
\]
O par\^{a}metro de ordem \(\Phi(t) e^{i\psi(t)} = N^{-1}\sum_j e^{i\theta_j}\), com \(\Phi \in [0,1]\), mede a coer\^{e}ncia de fase. Uma transi\c{c}\~{a}o de segunda ordem ocorre em \(K_c = 2/(\pi g(0))\): abaixo de \(K_c\), \(\Phi = 0\) (decoer\^{e}ncia total); acima, \(\Phi > 0\) (sincroniza\c{c}\~{a}o parcial ou total).

No DEQ\textsuperscript{2}, \(\Phi^{\text{interno}}\) de um ator mede a coer\^{e}ncia espontânea de seus subsistemas (institucional, econômico, diplomático, civil); \(\Phi^{\text{exportável}}\) mede sua capacidade de acoplar coer\^{e}ncia a subsistemas externos sem impor alinhamento ideológico.

\subsection{O funcional de estabilidade sistêmica}

O funcional central do DEQ\textsuperscript{2} \'{e} \citep{marcelino2026c}:
\begin{equation}\label{eq:ts-sys}
T_s^{\text{sys}} = \frac{s_{iii} \, z \, \delta \, (1 + \Phi \langle A \rangle)}{\sigma \, \varepsilon_{\text{dis}} \, (1 + \langle\Omega\rangle / \Phi)},
\end{equation}
onde:
\begin{itemize}\tightlist
\item \(s_{iii}\) \'{e} a legitimidade \'{e}tico-discursiva do ator;
\item \(z\) \'{e} a complexidade dial\'{e}tica m\'{e}dia do sistema;
\item \(\delta\) \'{e} o fator de desconto (horizonte temporal);
\item \(A\) \'{e} a antifragilidade \citep{taleb2012}: \(A = \partial^2\Pi/\partial\sigma^2\);
\item \(\Omega\) \'{e} a dissipa\c{c}\~{a}o involutiva (energia desperdi\c{c}ada em competi\c{c}\~{a}o autoprejudicial);
\item \(\sigma\) \'{e} a volatilidade estoc\'{a}stica;
\item \(\varepsilon_{\text{dis}}\) \'{e} o coeficiente de desengajamento moral \citep{bandura1999}.
\end{itemize}

O acoplamento assim\'{e}trico entre \(A\) e \(\Omega\) \'{e} a propriedade formal central: o valor (antifragilidade) requer transporte coerente para se propagar, portanto \(\Phi\) multiplica; o desperd\'{i}cio (involução) se autorreplica na aus\^{e}ncia de inibi\c{c}\~{a}o, portanto \(\Phi\) divide. Uma transi\c{c}\~{a}o de fase ocorre quando \(T_s^{\text{sys}} < 1\) por dura\c{c}\~{a}o \(\Delta t > \Delta t^*\).

\subsection{Os quatro atores da Parte III}

O DEQ\textsuperscript{2} classifica os quatro atores estruturais do sistema global 2024--2031 \citep{marcelino2026c}:

\begin{table}[h]
\centering
\caption{Tipologia DEQ\textsuperscript{2} dos quatro atores estruturais.}
\label{tab:atores}
{\small\begin{tabular}{llll}
\toprule
Ator & Tipo \(\Phi\) & Exportabilidade & Dire\c{c}\~{a}o \(T_s^{\text{sys}}\) \\
\midrule
Soliton (Musk) & \(\Phi\) degenerada (\(N=1\)) & por decoer\^{e}ncia for\c{c}ada & indefinida \\
Egem\^{o}n Decoerente (EUA) & \(\Phi \to 0\) & nenhuma & \(\to 0\) em 2027--2029 \\
R\'{i}gido Sincronizado (China) & \(\Phi\) imposta & exporta\c{c}\~{a}o de \(K\), n\~{a}o \(\Phi\) & \(\to 0\) catastr\'{o}fico \\
\textbf{Balan\c{c}a Qu\^{a}ntica (Brasil)} & \(\Phi\) espont\^{a}nea condicional & \textbf{\(\Phi\) real se A/B vencer} & \textbf{\(> 0\) se A/B; \(\to 0\) se C} \\
\bottomrule
\end{tabular}}
\end{table}

% ─── 3. O BRASIL COMO Φ-ENGINE ────────────────────────────────────────────────
\section{O Brasil como \(\Phi\)-engine: defini\c{c}\~{a}o formal e verifica\c{c}\~{a}o}

\begin{definition}[\(\Phi\)-\textit{engine} DEQ\textsuperscript{2}]
\label{def:phi-engine}
Um ator \(e\) \'{e} denominado \(\Phi\)-\textit{engine} se satisfaz simultaneamente:
\begin{enumerate}\tightlist
\item \(\Phi_e^{\text{interno}}\) moderado e \textbf{produzido espontaneamente} (n\~{a}o por coerc\c{c}\~{a}o hier\'{a}rquica);
\item \(\Phi_e^{\text{exportável}} > 0\) --- exist\^{e}ncia de \textit{interfaces de acoplamento bidirecional} com subsistemas globais heterog\^{e}neos;
\item \textbf{Assimetria \(A/\Omega\) favor\'{a}vel}: antifragilidade adaptativa, dissipa\c{c}\~{a}o involutiva abaixo da m\'{e}dia para sua faixa;
\item \(s_{iii,e} > s_{ii,e}\) relativamente (legitimidade excede capacidade material --- caracter\'{i}stica de \textit{bridge-builder});
\item \textbf{Pol\'{i}tica externa n\~{a}o-alinhada operativa}: o Estado Log\'{i}stico de Cervo como pol\'{i}tica efetiva.
\end{enumerate}
\end{definition}

\begin{proposition}\label{prop:brasil}
O Brasil de 2024--2026 satisfaz as cinco condi\c{c}\~{o}es da Defini\c{c}\~{a}o~\ref{def:phi-engine}, condicionalmente \`{a} continuidade operativa do Estado Log\'{i}stico após outubro de 2026.
\end{proposition}

\begin{proof}[Verifica\c{c}\~{a}o]
Condi\c{c}\~{a}o~(1): V-Dem registra U-turn democr\'{a}tico, com STF assertiva, independ\^{e}ncia de m\'{i}dia e neutralidade das For\c{c}as Armadas confirmada 2023--2025 \citep{vdem2026}. Coer\^{e}ncia espontânea, n\~{a}o coercitiva.

Condi\c{c}\~{a}o~(2): Presid\^{e}ncia BRICS 2025 (126 pontos de declara\c{c}\~{a}o, TFFF, BMG) e COP30 Bel\'{e}m (Mecanismo Bel\'{e}m, duas \textit{roadmaps}). Acoplamento bidirecional com Sul global (BRICS), Norte global (COP), e arena de conflito (Russia--Ucrânia via \textit{Group of Friends for Peace}).

Condi\c{c}\~{a}o~(3): resili\^{e}ncia institucional verificada (tentativa de golpe janeiro 2023, rejeitada; condena Bolsonaro setembro 2025). Dissipa\c{c}\~{a}o involutiva presente (polariza\c{c}\~{a}o afetiva, fiscal) mas abaixo do limiar de crise sistêmica.

Condi\c{c}\~{a}o~(4): Brand Finance Soft Power Index 2026: Brasil 49,2/100 (avan\c{c}o de 6 pp em dois anos). GDP PPP per capita médio-baixo entre os G20. \(s_{iii} > s_{ii}\) confirmado.

Condi\c{c}\~{a}o~(5): mantém rela\c{c}\~{o}es operativas com EUA, China, R\'{u}ssia, Ir\~{a} e Israel simultaneamente; participa do BRICS sem abandonar o \textit{track} OCDE; n\~{a}o aplicou san\c{c}\~{o}es unilaterais. Estado Log\'{i}stico operativo.
\end{proof}

% ─── 4. O ESTADO LOGÍSTICO ───────────────────────────────────────────────────
\section{O Estado Log\'{i}stico como realiza\c{c}\~{a}o pol\'{i}tica da \(\Phi\)-\textit{engine}}

Amado Luiz Cervo \citep{cervo2008} define o Estado Log\'{i}stico como paradigma em que o Estado n\~{a}o abre m\~{a}o da autonomia decisional mas a articula atrav\'{e}s de m\'{u}ltiplos vetores simult\^{a}neos: industrial, agr\'{i}cola, energ\'{e}tico, diplom\'{a}tico e ambiental. \'{E} distinto do Estado Normal (subordinado), do Estado Desenvolvimentista (autárquico) e do Estado Neoliberal (dependente).

A mapeamento DEQ\textsuperscript{2} formal do Estado Log\'{i}stico sobre o campo \(\Pi\) topogr\'{a}fico \'{e}:
\begin{equation}\label{eq:logistico}
\Pi^{\text{Log\'{i}stico}}(t+1) = \sum_{k} M_e^{(k)} \cdot Q_0^{(k)} \cdot (1 - \delta_R^{(k)} z_R^{(k)}),
\end{equation}
onde \(k\) indexa os vetores (industrial, agr\'{i}cola, energ\'{e}tico, diplom\'{a}tico, ambiental). \'{E} uma \textit{superposi\c{c}\~{a}o coerente} de m\'{u}ltiplos canais --- exatamente a estrutura que gera \(\Phi > 0\) num sistema multi-canal, em contraste com:

\begin{itemize}\tightlist
\item o Soliton, que maximiza um \'{u}nico canal (\(N=1\));
\item o Egem\^{o}n Decoerente, que perdeu coer\^{e}ncia entre seus canais (\(\Phi \to 0\));
\item o R\'{i}gido Sincronizado, que imp\~{o}e coer\^{e}ncia por coerc\c{c}\~{a}o (\(K \gg K_c\)).
\end{itemize}

\begin{remark}
O n\~{a}o-alinhamento ativo do Estado Log\'{i}stico n\~{a}o \'{e} neutralidade passiva (cf. n\~{a}o-alinhamento da Guerra Fria). \'{E} \textit{n\~{a}o-alinhamento generativo}: produz novos acoplamentos (TFFF, \textit{Group of Friends for Peace}, Mecanismo Bel\'{e}m) sem impor alinhamento sub-global. Esta \'{e} a propriedade operacional que distingue a \(\Phi\)-\textit{engine} da mera "pot\^{e}ncia m\'{e}dia".
\end{remark}

% ─── 5. EVIDÊNCIAS DE 2025 ───────────────────────────────────────────────────
\section{Tr\^{e}s evid\^{e}ncias estruturais de 2025}

\subsection{BRICS Rio de Janeiro --- Julho de 2025}

O Brasil presidiu o 17\textdegree{} summit BRICS (Rio de Janeiro, 6--7 de julho de 2025). A declara\c{c}\~{a}o final de 126 pontos incluiu: o \textit{Tropical Forest Forever Fund} (TFFF) com meta de capitaliza\c{c}\~{a}o de \$10B até 2030; o \textit{BRICS Multilateral Guarantees} (BMG) para infraestrutura no Sul global; a declara\c{c}\~{a}o sobre governan\c{c}a global da IA. Foram realizadas mais de 200 reuniões preparatórias.

\textbf{Leitura DEQ\textsuperscript{2}:} o BRICS-2025 mede \(\Phi^{\text{Sul global}}\) --- capacidade de síntese entre interesses divergentes (China--Índia, Rússia--Iran, Brasil--mediador). A aprova\c{c}\~{a}o de instrumentos financeiros n\~{a}o-triviais por consenso de 10 membros com interesses freq\"{u}entemente opostos constitui evidência emp\'{i}rica de \(\Phi^{\text{export\'{a}vel}} > 0\).

\subsection{COP30 Bel\'{e}m --- Novembro de 2025}

A COP30 (Bel\'{e}m, 10--22 de novembro de 2025) foi presidida pelo Brasil at\'{e} novembro de 2026. Os tr\^{e}s pilares da presid\^{e}ncia brasileira produziram: mobiliza\c{c}\~{a}o de \$1,3T/ano at\'{e} 2035 para clima; triplicamento do financiamento de adapta\c{c}\~{a}o; o \textit{Mecanismo Bel\'{e}m para Transi\c{c}\~{a}o Justa Global}; duas \textit{roadmaps} lan\c{c}adas pela presid\^{e}ncia apesar da aus\^{e}ncia de consenso formal (\(\approx 80\) pa\'{i}ses pró vs.\ \(\approx 80\) contra linguagem explícita).

\textbf{Leitura DEQ\textsuperscript{2}:} a COP30 mede \(\Phi^{\text{Norte--Sul}}\) --- capacidade de mediar entre altas emiss\~{o}es e adapta\c{c}\~{a}o. O lan\c{c}amento de dois documentos de \textit{roadmap} apesar da divis\~{a}o formal constitui uma jogada de lideran\c{c}a n\~{a}o-hegemônica: produz estrutura sem impor alinhamento.

\subsection{A condena Bolsonaro --- Setembro e Novembro de 2025}

Em 11 de setembro de 2025, o STF condenou Jair Bolsonaro a 27 anos e 3 meses por: tentativa de golpe de Estado (8 de janeiro de 2023); participa\c{c}\~{a}o em organiza\c{c}\~{a}o criminosa armada; plano para assassinar Lula, Alckmin e Moraes; aboli\c{c}\~{a}o violenta da ordem democr\'{a}tica. Em 7 de novembro de 2025, o painel do STF confirmou a condena (4--1, dissid\^{e}ncia do juiz Fux). \'{E} a primeira condena na hist\'{o}ria brasileira de um ex-presidente por tentativa de golpe.

\textbf{Leitura DEQ\textsuperscript{2}:} este evento mede a \textit{resist\^{e}ncia cr\'{i}tica} \(R_c\) do sistema brasileiro \citep{marcelino2026a}: a capacidade do subsistema judicial de se opor \`{a} trajet\'{o}ria de um ator dominante sob press\~{a}o m\'{a}xima. \(R_c > 0\) efetivo \'{e} condi\c{c}\~{a}o necess\'{a}ria de \(\Phi\) espont\^{a}nea. Nenhum dos outros tr\^{e}s atores do Quadro~\ref{tab:atores} exibe esta propriedade: EUA erodem os checks institucionais (Cap.~7 DEQ\textsuperscript{2}), China n\~{a}o os admite por design (Cap.~8), Soliton n\~{a}o est\'{a} sujeito a eles.

% ─── 6. FRAGILIDADES REAIS ───────────────────────────────────────────────────
\section{As tr\^{e}s fragilidades: honestidade anal\'{i}tica}

O argumento perderia rigor se omitisse as fragilidades reais. A \(\Phi\)-\textit{engine} brasileira n\~{a}o \'{e} autom\'{a}tica.

\subsection{Polariza\c{c}\~{a}o e aprova\c{c}\~{a}o em queda}

Aprova\c{c}\~{a}o Lula em mar\c{c}o de 2026 (Genial/Quaest, 2.004 eleitores): 44\% aprova, 51\% desaprova --- pior gap desde julho de 2025. A coligação de governo atinge m\'{i}nimo hist\'{o}rico de 30 anos em fidelidade parlamentar. A polariza\c{c}\~{a}o afetiva \'{e} o ru\'{i}do end\'{o}geno que erode \(\Phi^{\text{interno}}\). O Brasil n\~{a}o est\'{a} imune ao fen\^{o}meno global --- apenas disp\~{o}e de mecanismos \(R_c\) ainda operantes para contê-lo.

\subsection{Quadro fiscal e Selic}

Crescimento 2025: +2,3\%. Previs\~{a}o 2026: +1,7\%. Selic de 10,5\% (setembro de 2024) a 15\% (junho de 2025); afrouxamento esperado para 9,75--11,50\% no final de 2026. Inflação IPCA 2025: 4,44\% (dentro da meta 1,5--4,5\%). D\'{e}ficit federal 2026: \(\approx 3\%\) PIB; d\'{i}vida p\'{u}blica \(\approx 75\%\) PIB.

O desequil\'{i}brio fiscal aumenta \(\sigma\) ex\'{o}gena: volatilidade monet\'{a}ria, press\~{a}o cambial, vulnerabilidade a choques externos. A combina\c{c}\~{a}o \textit{crescimento desacelerado + Selic alta + d\'{i}vida 75\%} reduz a capacidade do governo de sustentar financeiramente a agenda do Estado Log\'{i}stico (TFFF, media\c{c}\~{a}o, COP).

\subsection{Elei\c{c}\~{o}es de 4 de outubro de 2026}

Bolsonaro inelegível (condenado). Seu filho \textbf{Fl\'{a}vio Bolsonaro}, candidato do Partido Liberal. Sondagem AtlasIntel/Bloomberg de abril de 2026: Fl\'{a}vio 47,6\% vs.\ Lula 46,6\% num segundo turno simulado --- primeira vez em que Fl\'{a}vio est\'{a} numericamente \`{a} frente (margem dentro do erro estat\'{i}stico). Primeiro turno: Lula 45,9\% vs.\ Fl\'{a}vio 40,1\%. Tend\^{e}ncia dos meses anteriores: Lula em queda de +12 pp (dezembro de 2025) para empate estat\'{i}stico em abril de 2026.

% ─── 7. CENÁRIOS ELEITORAIS ──────────────────────────────────────────────────
\section{Tr\^{e}s cen\'{a}rios e suas consequ\^{e}ncias para \(T_s^{\text{sys}}\)}

\begin{table}[h]
\centering
\caption{Cen\'{a}rios eleitorais de outubro de 2026 e consequ\^{e}ncias DEQ\textsuperscript{2}.}
\label{tab:cenarios}
{\small\begin{tabular}{p{3.2cm}p{2.8cm}p{5.5cm}}
\toprule
Cen\'{a}rio & \(P\) subjetiva (abril 2026) & Consequ\^{e}ncia para \(\Phi\)-\textit{engine} \\
\midrule
\textbf{A} Lula re-eleito & 0,45--0,55 & \(\Phi\)-\textit{engine} continua sob press\~{a}o fiscal mas operativa; TFFF e BMG avan\c{c}am \\
\textbf{B} Sucessor de esquerda (ex.: Haddad) & 0,10--0,15 & \(\Phi\)-\textit{engine} continua com poss\'{i}vel reequil\'{i}brio fiscal; trajet\'{o}ria análoga a A \\
\textbf{C} Fl\'{a}vio Bolsonaro ou direita radical & 0,30--0,40 & Interrup\c{c}\~{a}o do Estado Log\'{i}stico; trajet\'{o}ria rumo a Egem\^{o}n Decoerente regional; \(T_s^{\text{sys}} \to 0\) em 24 meses \\
\bottomrule
\end{tabular}}
\end{table}

A tese central do DEQ\textsuperscript{2} para este ponto de bifurca\c{c}\~{a}o:

\begin{itemize}\tightlist
\item \textbf{Se A ou B:} consolidação do Brasil como \(\Phi\)-\textit{engine} global at\'{e} 2028--2031, com \(T_s^{\text{sys},\text{BRA}}\) significativamente superior a EUA e China segundo o funcional \eqref{eq:ts-sys}.

\item \textbf{Se C:} solitoniza\c{c}\~{a}o discursiva do Brasil --- transi\c{c}\~{a}o do perfil Dial\'{o}gico \((7, 6, 8)\) para o Egem\^{o}nico-Decoerente \((8, 9, 2)\), recalcando a trajet\'{o}ria norte-americana --- com perda parcial de \(s_{iii}\) global em 24 meses.
\end{itemize}

\begin{remark}
A probabilidade de que o sistema global 2028--2031 entre em sua janela de transi\c{c}\~{a}o de fase com pelo menos um produtor de \(\Phi\) export\'{a}vel \'{e} \(P \approx 0,55\text{--}0,70\), condicional ao resultado eleitoral de outubro de 2026. Esta \'{e} a \'{u}nica vari\'{a}vel do sistema cujo valor ainda n\~{a}o foi fixado pelos outros tr\^{e}s atores da Tabela~\ref{tab:atores}.
\end{remark}

% ─── 8. PREDIÇÕES FALSIFICÁVEIS ──────────────────────────────────────────────
\section{Quatro predi\c{c}\~{o}es datadas pr\'{e}-registradas}

Pr\'{e}-registro p\'{u}blico: \texttt{github.com/bilanciere-quantistico/predizioni-deq2} (predi\c{c}\~{o}es B1--B4 na tabela do reposit\'{o}rio).

\subsection{B1 --- Resultado da elei\c{c}\~{a}o de 4 de outubro de 2026}

\textbf{Predi\c{c}\~{a}o:} probabilidade de vit\'{o}ria de Lula ou sucessor de esquerda \(P = 0,55\text{--}0,70\) em abril de 2026; estima\c{c}\~{a}o a ser revisada em 30 de junho e 30 de setembro de 2026 com base em m\'{e}dia ponderada de Genial/Quaest, AtlasIntel e Datafolha.

\textbf{Condi\c{c}\~{a}o de falsifica\c{c}\~{a}o:} vit\'{o}ria de candidato de esquerda com margem \(< 1\%\) no primeiro turno e \(< 2\%\) no segundo torna a predi\c{c}\~{a}o indeterminada (resultado condicional, n\~{a}o refutado).

\textbf{Prazo:} 4 de outubro de 2026. Fonte: TSE (resultado oficial).

\subsection{B2 --- TFFF e BRICS BMG operativos (condicional a B1 = A/B)}

\textbf{Predi\c{c}\~{a}o:} se Lula ou sucessor de esquerda eleito, o TFFF recebe \(>\,\$5\)B comprometidos at\'{e} o final de 2027 e o BMG opera\c{c}\~{o}es se iniciam antes de Q2 2027.

\textbf{Condi\c{c}\~{a}o de falsifica\c{c}\~{a}o:} TFFF capitaliza\c{c}\~{a}o \(< \$1\)B at\'{e} final de 2027 OU BMG sem opera\c{c}\~{o}es iniciadas at\'{e} Q2 2027.

\textbf{Prazo:} Q2 2027. Fonte: comunicados BRICS, Banco Mundial.

\subsection{B3 --- V-Dem Bra\-sil em alta (condicional a B1 = A/B)}

\textbf{Predi\c{c}\~{a}o:} \'{i}ndice V-Dem Liberal Democracy do Brasil ultrapassa 0,65 no relat\'{o}rio de 2028 (referente ao ano de 2027).

\textbf{Condi\c{c}\~{a}o de falsifica\c{c}\~{a}o:} \'{i}ndice V-Dem 2027 abaixo de 0,60 ou deteriora\c{c}\~{a}o de mais de 0,03 pontos em rela\c{c}\~{a}o a 2025.

\textbf{Prazo:} relat\'{o}rio V-Dem publicado em abril de 2028. Fonte: \texttt{v-dem.net}.

\subsection{B4 --- Deteriora\c{c}\~{a}o estrutural em caso C}

\textbf{Predi\c{c}\~{a}o:} se Fl\'{a}vio Bolsonaro ou candidato de direita radical vencer em outubro de 2026, o Brasil registra: deteriora\c{c}\~{a}o V-Dem \(> 0,05\) pontos em 24 meses; sa\'{i}da ou passiviza\c{c}\~{a}o no \textit{Group of Friends for Peace}; participa\c{c}\~{a}o BRICS transformada de ativa para formal-passiva.

\textbf{Condi\c{c}\~{a}o de falsifica\c{c}\~{a}o:} se C vencer mas V-Dem permanecer est\'{a}vel e o Estado Log\'{i}stico continuar operativo, a teoria DEQ\textsuperscript{2} requer revis\~{a}o do peso relativo da vari\'{a}vel eleitoral sobre a trajet\'{o}ria de \(\Phi\).

\textbf{Prazo:} outubro de 2028. Fontes: V-Dem, BRICS, ONU.

% ─── 9. CONCLUSÃO ────────────────────────────────────────────────────────────
\section{Conclus\~{a}o}

O Brasil de 2025--2026 constitui o \'{u}nico ator do sistema global com a estrutura formal de um \(\Phi\)-\textit{engine} no sentido DEQ\textsuperscript{2}: coer\^{e}ncia espont\^{a}nea, exportabilidade verificada nas duas presi\-d\^{e}ncias de 2025, assimetria \(A/\Omega\) favor\'{a}vel confirmada pela condena Bolsonaro como prova de \(R_c\) efetivo, e \(s_{iii} > s_{ii}\) documentado. O Estado Log\'{i}stico de Cervo n\~{a}o \'{e} retórica diplom\'{a}tica: \'{e} a realiza\c{c}\~{a}o pol\'{i}tica de uma estrutura matem\'{a}tica que o modelo DEQ\textsuperscript{2} formaliza como superposi\c{c}\~{a}o coerente de m\'{u}ltiplos canais.

Mas a fun\c{c}\~{a}o \'{e} condicional. A elei\c{c}\~{a}o de 4 de outubro de 2026 \'{e} o ponto de bifurca\c{c}\~{a}o expl\'{i}cito da teoria: com data conhecida, mecanismo formalizado, predi\c{c}\~{o}es pr\'{e}-registradas e condi\c{c}\~{o}es de falsifica\c{c}\~{a}o expl\'{i}citas. Se vencer A ou B: o Brasil entrar\'{a} na janela 2028--2031 como o \'{u}nico produtor de \(\Phi\) export\'{a}vel do sistema, com consequ\^{e}ncias estruturais para a governan\c{c}a global do clima, das finan\c{c}as do Sul e da mediação de conflitos. Se vencer C: a teoria prevê solitoniza\c{c}\~{a}o discursiva e deteriora\c{c}\~{a}o institucional mensur\'{a}veis em 24 meses.

A singularidade do caso brasileiro no DEQ\textsuperscript{2} n\~{a}o reside na pot\^{e}ncia material --- que \'{e} m\'{e}dia. Reside na combina\c{c}\~{a}o de todos os demais fatores: o Estado Log\'{i}stico como pr\'{a}tica, a resili\^{e}ncia democr\'{a}tica como prova, a dupla presid\^{e}ncia multilateral como demonstra\c{c}\~{a}o. Nenhum outro ator do sistema global 2024--2026 exibe esta combina\c{c}\~{a}o. Esta \'{e} a afirma\c{c}\~{a}o central do artigo --- e \'{e} falsific\'{a}vel.

\section*{Refer\^{e}ncias}

\begin{thebibliography}{99}

\bibitem[Bandura(1999)]{bandura1999}
Bandura, A. (1999). Moral disengagement in the perpetration of inhumanities. \textit{Personality and Social Psychology Review}, 3(3), 193--209.

\bibitem[Brand Finance(2026)]{brandfinance2026}
Brand Finance. (2026). \textit{Brand Finance Soft Power Index 2026}. Brand Finance, London.

\bibitem[BRICS Brasil(2025)]{brics2025}
BRICS Brasil. (2025). \textit{XVII BRICS Summit Leaders' Declaration}. Rio de Janeiro, 6--7 de julho de 2025.

\bibitem[Carbon Brief(2025)]{carbonbrief2025}
Carbon Brief. (2025). \textit{COP30: Key outcomes from the Belém climate summit}. Carbon Brief, London.

\bibitem[Cervo(2008)]{cervo2008}
Cervo, A. L. (2008). \textit{Inser\c{c}\~{a}o internacional: forma\c{c}\~{a}o dos conceitos brasileiros}. S\~{a}o Paulo: Saraiva.

\bibitem[Cervo e Lessa(2014)]{cervoe2014}
Cervo, A. L., \& Lessa, A. C. (2014). O declínio: inser\c{c}\~{a}o internacional do Brasil (2011--2014). \textit{Revista Brasileira de Pol\'{i}tica Internacional}, 57(2), 133--151.

\bibitem[Furtado(1974)]{furtado1974}
Furtado, C. (1974). \textit{O mito do desenvolvimento econômico}. Rio de Janeiro: Paz e Terra.

\bibitem[Kuramoto(1984)]{kuramoto1984}
Kuramoto, Y. (1984). \textit{Chemical Oscillations, Waves, and Turbulence}. Springer, Berlin.

\bibitem[Lessa et al.(2013)]{lessa2013}
Lessa, A. C., Couto, L. F., \& Farias, R. S. (2013). Pol\'{i}tica exterior do Brasil no contexto da globaliza\c{c}\~{a}o. \textit{Revista Brasileira de Pol\'{i}tica Internacional}, 56(1), 73--91.

\bibitem[Marcelino(2026a)]{marcelino2026a}
Marcelino, L. F. (2026a). \textit{A Topografia do Discurso Emancipat\'{o}rio}. KDP, Mil\~{a}o.

\bibitem[Marcelino(2026b)]{marcelino2026b}
Marcelino, L. F. (2026b). \textit{A Geometria da Hegemonia}. KDP, Mil\~{a}o.

\bibitem[Marcelino(2026c)]{marcelino2026c}
Marcelino, L. F. (2026c). \textit{A Balan\c{c}a Qu\^{a}ntica: Geometria da Hegemonia e Transi\c{c}\~{a}o de Fase 2028--2031}. KDP, Mil\~{a}o.

\bibitem[Marcelino e Teixeira(2026)]{marcelino2026ct}
Marcelino, L. F., \& Teixeira, E. (2026). \textit{A Constante Marcelino}. KDP, Mil\~{a}o.

\bibitem[Munich Security(2025)]{munich2025}
Munich Security Conference. (2025). \textit{Munich Security Report 2025: Lose-Lose?} Munich: MSC.

\bibitem[Pinheiro(2004)]{pinheiro2004}
Pinheiro, L. (2004). Pol\'{i}tica externa brasileira. \textit{Puc-Rio}.

\bibitem[Saraiva(2010)]{saraiva2010}
Saraiva, M. G. (2010). As estratégias de coopera\c{c}\~{a}o Sul-Sul nos marcos da pol\'{i}tica externa brasileira de Lula a Dilma Rousseff. \textit{Revista Brasileira de Pol\'{i}tica Internacional}, 53(2), 151--168.

\bibitem[Taleb(2012)]{taleb2012}
Taleb, N. N. (2012). \textit{Antifragile: Things That Gain from Disorder}. Random House, New York.

\bibitem[V-Dem(2026)]{vdem2026}
V-Dem Institute. (2026). \textit{Democracy Report 2026: Democracy Winning and Losing at the Ballot}. Gothenburg: University of Gothenburg.

\bibitem[Vigevani e Cepaluni(2007)]{vigevani2007}
Vigevani, T., \& Cepaluni, G. (2007). A pol\'{i}tica externa de Lula da Silva: a estrat\'{e}gia da autonomia pela diversifica\c{c}\~{a}o. \textit{Contexto Internacional}, 29(2), 273--335.

\bibitem[Vizentini(2005)]{vizentini2005}
Vizentini, P. G. F. (2005). De FHC a Lula: uma d\'{e}cada de pol\'{i}tica externa (1995--2005). \textit{Civitas -- Revista de Ci\^{e}ncias Sociais}, 5(2), 381--397.

\end{thebibliography}

\end{document}
